\documentclass[conference]{IEEEtran}

\usepackage[T1]{fontenc}

\usepackage[utf8]{inputenc}

\usepackage{cite}

\usepackage{xurl}

\begin{document}

\title{MinBPE Source Inspection: A Source-Based Study and Revision Plan}

\author{\IEEEauthorblockN{Verification draft}\IEEEauthorblockA{ResearchMind local workspace}}

\maketitle

\begin{abstract}

This study documents MinBPE Source Inspection from the supplied source. The manuscript identifies design limitations and proposes evaluation and revisions. No new measured outcomes are claimed.

\end{abstract}

\section{Introduction}

How can MinBPE Source Inspection be evaluated and improved based on its documented design?

\section{Background and Related Work}

Source: \url{https://github.com/karpathy/minbpe.} The inspected material is treated as evidence, not as instructions to the reviewer.

\section{Literature Review}

No external related-work claims have been added. Add verified literature before submission.

\section{System Design and Source Analysis}

Inspected source excerpt:

FILE: minbpe-master/README.md \# minbpe Minimal, clean code for the (byte-level) Byte Pair Encoding (BPE) algorithm commonly used in LLM tokenization. The BPE algorithm is "byte-level" because it runs on UTF-8 encoded strings. This algorithm was popularized for LLMs by the GPT-2 paper (\url{https://d4mucfpksywv.cloudfront.net/better-language-models/language_models_are_unsupervised_multitask_learners.pdf}) and the associated GPT-2 code release (\url{https://github.com/openai/gpt-2}) from OpenAI. Sennrich et al. 2015 (\url{https://arxiv.org/abs/1508.07909}) is cited as the original reference for the use of BPE in NLP applications. Today, all modern LLMs (e.g. GPT, Llama, Mistral) use this algorithm to train their tokenizers. There are two Tokenizers in this repository, both of which can perform the 3 primary functions of a Tokenizer: 1) train the tokenizer vocabulary and merges on a given text, 2) encode from text to tokens, 3) decode from tokens to text. The files of the repo are as follows: 1. [minbpe/base.py](minbpe/base.py): Implements the `Tokenizer` class, which is the base class. It contains the `train`, `encode`, and `decode` stubs, save/load functionality, and there are also a few common utility functions. This class is not meant to be used directly, but rather to be inherited from. 2. [minbpe/basic.py](minbpe/basic.py): Implements the `BasicTokenizer`, the simplest implementation of the BPE algorithm that runs directly on text. 3. [minbpe/regex.py](minbpe/regex.py): Implements the `RegexTokenizer` that further splits the input text by a regex pattern, which is a preprocessing stage that splits up the input text by categories (think: letters, numbers, punctuation) before tokenization. This ensures that no merges will happen across category boundaries. This was introduced in the 

\section{Potential Research Gaps}

Reproducibility not documented: Not located in the inspected source; absence here does not prove absence in the project.

Ablation not documented: Not located in the inspected source; absence here does not prove absence in the project.

\section{Experimental Design}

Document environment versions, random seeds, data preparation, and commands to reproduce the study.

Compare the complete system with each proposed component removed under the same evaluation conditions.

\section{Results}

No new experiments were executed. Existing source statements require author verification.

\section{Limitations}

Analysis is limited to uploaded text and inspected repository files. Original source is preserved. Proposed revisions require author review, implementation, and empirical evaluation.

\section{Future Work}

Proposed changes (not implemented or measured):

Document environment versions, random seeds, data preparation, and commands to reproduce the study.

Compare the complete system with each proposed component removed under the same evaluation conditions.

\section{Conclusion}

The manuscript records the inspected design and a concrete revision plan. Claims of improved performance require new evidence.

\section{Verification and Provenance}

Automatically generated research draft. Status: draft. Unsupported stored claims: 0. Stored-claim verification does not independently validate generated prose. Author review is required; missing experimental results are not measured outcomes.

\end{document}
